% 第11章 均值回归策略当前局限性摘要
\section{均值回归策略当前局限性摘要}\label{sec:limitations}

基于第 \ref{sec:survey} 章分方法族综述与第 \ref{sec:audit} 章污染审计，本摘要归纳均值回归策略从“理论事实”走向“可持续实践”的九大局限。每条按“本质问题 $\to$ 证据 $\to$ 实践含义”组织，并附置信度。

\begin{enumerate}
  \item \textbf{利润的时间衰减（置信度：高）。} 本质问题：均值回归的超额收益并非时不变常量，而是随市场参与者竞争而衰减。证据：Gatev 等（2006）\cite{gatev2006} 在 1962--2002 年样本中年化超额约 11\%，而 Do \& Faff（2012）\cite{dofaff2012} 将样本延伸至 2009 年后，发现利润系统性下降，后期几乎被交易成本吞噬。实践含义：任何“复制历史基准”的结论都必须附样本区间与衰减曲线，不能直接外推。

  \item \textbf{交易成本的侵蚀性（置信度：高）。} 本质问题：均值回归信号天然高频（持有期短、换手高），每一次进出都支付价差与冲击成本。证据：Martin \& Sch\"oneborn（2011）\cite{martin2011} 的理论与数值研究表明“均值回归可盈利，但成本决定成败”；回测中零成本与真实成本两档收益可差一个量级。实践含义：成本口径必须与策略频率匹配；冲击成本模型（如绝对参与率）应纳入优化而非事后扣减。

  \item \textbf{机制归因的混乱（置信度：高）。} 本质问题：“反转利润”可能来自过度反应、流动性提供、交叉自相关（lead--lag）等多种机制，机制不同则可持续性不同。证据：Lo \& MacKinlay（1990）\cite{lomackinlay1990} 证明聚合层面的反转利润主要来自交叉自相关而非个股过度反应；短期反转（Lehmann 1990 \cite{lehmann1990}）利润部分来自买卖价差。实践含义：策略报告必须交代利润的机制来源，否则无法判断它能否在交易成本变化或市场结构变化下存活。

  \item \textbf{参数非平稳与估计误差（置信度：中--高）。} 本质问题：OU 过程的回归强度 $\theta$、长期均值 $\mu$、半衰期 $H=\ln 2/\theta$ 均随样本漂移，点估计误差极大。证据：Elliott 等（2005）\cite{elliott2005} 用状态空间在线估计即是应对非平稳的尝试，但其状态先验仍依赖建模者选择。实践含义：上线前必须做参数滚动估计的敏感性检验，并以区间/状态估计而非点估计驱动信号。

  \item \textbf{容量限制（置信度：中）。} 本质问题：统计套利与配对交易的收益来自有限个价差机会，资金规模增大后收益被冲击成本与拥挤交易稀释。证据：Avellaneda \& Lee（2010）\cite{avellaneda2010} 明确讨论策略容量与衰减；Gatev 距离法的配对池同样有限。实践含义：回测报告必须给出容量估计与参与率上限，而不是只报百分比收益。

  \item \textbf{评测污染的系统性存在（置信度：高）。} 本质问题：均值回归文献普遍存在数据窥探、幸存者偏差与 look-ahead 风险，且 2010 年前的研究几乎未做多重检验校正。证据：White（2000）\cite{white2000}、Sullivan 等（1999）\cite{sullivan1999}、Harvey 等（2016）\cite{harveylz2016}、Bailey 等 \cite{bailey2014pbo,bailey2014dsr} 构成完整的审计工具链，但经典反转文献极少使用。实践含义：在引用任何“高收益声明”前，先检查其是否报告试验次数与 OOS 切分；没有报告的，默认按多重检验风险处理。

  \item \textbf{机器学习方法的证据不足（置信度：低--中）。} 本质问题：深度/强化学习方法在 arXiv 上宣称优于经典基线，但独立复现、公开基准与多重检验披露普遍缺位。证据：复核后（见第 \ref{sec:checklist} 章）确认，14 篇预印本中仅个别工作报告了严格的样本外研究——如 R34（arXiv:2106.04028）对 500 只美股 19 年日频做样本外检验，报告年化夏普约 4 并检验交易成本敏感性；但多数预印本未披露试验次数、完整成本模型或独立复现代码，加密与新兴市场样本的稳健性证据更弱。实践含义：在独立复现通过之前，ML 增量应视为“设定红利”而非“模型红利”。

  \item \textbf{市场微观结构演化（置信度：中）。} 本质问题：高频做市与算法交易的普及改变了反转信号的形成与衰减速度，历史反转动力学未必延续。证据：微结构层面的反转交易（arXiv:2608.00885 \cite{arxiv_microstructure}）与早期周/月反转属于不同时间尺度，其驱动机制（订单流不均衡、存货效应）差异显著。实践含义：必须区分“微观结构反转”与“行为金融反转”，两者的数据频率、容量与成本结构完全不同。

  \item \textbf{Regime 依赖（置信度：中--高）。} 本质问题：反转与动量在同一资产上并存，何者占优取决于时间尺度与市场状态；均值回归策略在趋势市（如商品单边行情）中持续回撤。证据：Moskowitz 等（2012）\cite{mom2012} 与 Erb \& Harvey（2006）\cite{erbharvey2006} 显示同批资产上中长周期动量与商品展期反转并存；反转策略在低波动震荡环境表现更佳。实践含义：均值回归仓位应与市场状态（Regime）判定耦合，并在反转失效时设置硬性退出，而非无条件“逢低买入”。
\end{enumerate}

\noindent\textbf{综合判断}：均值回归的“理论事实”稳健（高置信度），但“可交易性”受成本、容量、机制与评测污染的叠加制约。当前文献给出的最优实践是——用 OU/协整框架做信号构造、用严格成本与多重检验做验证、用 Regime 判定做仓位管理、用跨市场 walk-forward 做泛化确认。任何只依赖单一样本、单一区间的均值回归“高收益”结论，都应按第 \ref{sec:reproduce} 章路线重验。
